\chapter{Introducción}

\section{Objetivo de la práctica}

El objetivo principal de esta práctica es configurar y utilizar \textbf{correo electrónico seguro} mediante el estándar \textbf{OpenPGP}, incluyendo:

\begin{itemize}[label=\textbullet]
    \item Generación de un par de claves OpenPGP personal en Mozilla Thunderbird
    \item Intercambio de claves públicas con la compañera (Wafa Azdad Triki)
    \item Envío y recepción de mensajes firmados y cifrados con OpenPGP
    \item Gestión de la confianza de claves (Web of Trust)
    \item Decodificación manual de mensajes cifrados con OpenPGP usando la utilidad \texttt{gpg} de GnuPG
\end{itemize}

\section{Conceptos fundamentales}

\subsection{OpenPGP y el modelo Web of Trust}

\textbf{OpenPGP} es un estándar criptográfico (RFC 4880 / RFC 9580) para cifrado y firma digital basado en clave pública. A diferencia de S/MIME, que confía en una jerarquía de Autoridades de Certificación (CA), OpenPGP utiliza un modelo denominado \textbf{Web of Trust} (red de confianza):

\begin{itemize}[label=\textbullet]
    \item No existe una CA central: cada usuario genera su propio par de claves.
    \item La confianza se establece entre individuos: otros usuarios firman tu clave pública para atestiguarla.
    \item Los niveles de confianza son configurables: desconocida, marginal, completa, o absoluta (\textit{ultimate}).
\end{itemize}

\subsection{Estructura de una clave OpenPGP}

Un par de claves OpenPGP típico en Thunderbird contiene:

\begin{table}[H]
    \centering
    \caption{Componentes de un par de claves OpenPGP}
    \begin{tabularx}{\textwidth}{l|l|X}
        \toprule
        \textbf{Componente} & \textbf{Uso} & \textbf{Descripción} \\
        \midrule
        Clave maestra (SC) & Sign + Certify & Firma mensajes y certifica otras claves \\
        Subclave (E) & Encrypt & Se usa exclusivamente para cifrar \\
        \bottomrule
    \end{tabularx}
\end{table}

En esta práctica se observa que la clave de mi compañera, ``Wafa'' tiene:
\begin{itemize}[label=\textbullet]
    \item Clave maestra \texttt{0x62C2D55D111C2F3C} (SC -- firma y certificación)
    \item Subclave \texttt{0xE1B59CED1F89D16C} (E -- cifrado)
\end{itemize}
Y nuestra clave tiene ID \texttt{0x4B0E6E00A737A803}, con subclave de descifrado \texttt{0xB809E1B2EC24D298}.

\subsection{Servicios de seguridad de OpenPGP}

\begin{table}[H]
    \centering
    \caption{Servicios de seguridad proporcionados por OpenPGP}
    \begin{tabularx}{\textwidth}{l|l|X}
        \toprule
        \textbf{Servicio} & \textbf{Mecanismo} & \textbf{Descripción} \\
        \midrule
        Confidencialidad & Cifrado híbrido & Cifrado simétrico de sesión, protegido con clave pública del destinatario \\
        Autenticación & Firma digital & Se verifica la identidad del emisor mediante su clave pública \\
        Integridad & Firma digital & Se detecta cualquier alteración del mensaje \\
        \bottomrule
    \end{tabularx}
\end{table}

\subsection{Formatos de ficheros}

\begin{table}[H]
    \centering
    \caption{Formatos de ficheros utilizados en esta práctica}
    \begin{tabularx}{\textwidth}{l|l|X}
        \toprule
        \textbf{Extensión} & \textbf{Formato} & \textbf{Contenido} \\
        \midrule
        .asc & ASCII Armored (Base64) & Clave pública o privada OpenPGP exportada \\
        .eml & Texto & Mensaje de correo completo (cabeceras + cuerpo MIME) \\
        .mime & Texto & Salida MIME del mensaje descifrado por gpg \\
        .gpg & Binario & Clave OpenPGP en formato binario (tras \texttt{--dearmor}) \\
        \bottomrule
    \end{tabularx}
\end{table}

\section{Herramientas utilizadas}

\begin{itemize}[label=\textbullet]
    \item \textbf{Mozilla Thunderbird} con soporte nativo OpenPGP: generación de claves, envío y recepción de mensajes cifrados y firmados.
    \item \textbf{GnuPG (gpg)}: herramienta de línea de comandos para importar claves, gestionar el llavero y descifrar/verificar mensajes manualmente incluida en la mayoría de distribuciones Linux.
    \item \textbf{ripmime}: utilidad para extraer adjuntos de mensajes MIME, usada cuando el mensaje descifrado contiene contenido codificado en Base64.
\end{itemize}

